\chapter{Algebre di Boole}\label{cap:boole}

In questo capitolo chiudiamo il cerchio. Nel Capitolo~\ref{cap:anelli} abbiamo incontrato gli \textbf{anelli booleani}, come
$(\mcP(s), \Delta, \cap)$; nel Capitolo~\ref{cap:reticoli} i reticoli distributivi e complementati, come $(\mcP(s), \subseteq)$. Vediamo che si tratta dello stesso
oggetto visto in tre modi: \textbf{reticolo booleano}, \textbf{algebra di Boole}, \textbf{anello booleano}. Il teorema di Stone dirà poi che,
nel caso finito, l'unico esempio è proprio l'insieme delle parti, e le funzioni caratteristiche lo trasformeranno nelle stringhe di bit
con le operazioni \texttt{AND}, \texttt{OR}, \texttt{XOR} e \texttt{NOT} che si usano in informatica.

% ──────────────────────────────────────────────────────────────
\section{Reticoli booleani e algebre di Boole}\label{sec:algebre-boole}

\dfn[reticolo-booleano]{Reticolo booleano}{
    Un reticolo distributivo e complementato si dice \textbf{reticolo booleano}. Indichiamo con $0$ il suo minimo, con $1$ il suo massimo e con
    $x'$ il complemento di $x$, che è unico per il Teorema~\ref{teo:complemento-unico}.
}

\dfn[algebra-boole]{Algebra di Boole}{
    Una struttura $(s, \vee, \wedge, 0, 1, {}')$, dove $\vee$ e $\wedge$ sono operazioni interne su $s$, $x \mapsto x'$ è una funzione da $s$ in $s$ e $0, 1 \in s$, si dice
    \textbf{algebra di Boole} se:
    \begin{enumerate}
        \item $(s, \wedge, 1)$ e $(s, \vee, 0)$ sono monoidi commutativi in cui ogni elemento è idempotente ($x \wedge x = x = x \vee x$);
        \item vale l'assorbimento: $x \vee (x \wedge y) = x$ e $x \wedge (x \vee y) = x$ per ogni $x, y$;
        \item $\wedge$ è distributiva rispetto a $\vee$ e $\vee$ è distributiva rispetto a $\wedge$;
        \item $x \vee x' = 1$ e $x \wedge x' = 0$ per ogni $x$.
    \end{enumerate}
}

\begin{esempio}{Due algebre di Boole}
    I due esempi principali:
    \begin{center}
    \renewcommand{\arraystretch}{1.3}
    \begin{tabular}{l|c|c|c|c|c|c}
        & ordine & $\vee$ & $\wedge$ & $0$ & $1$ & $x'$ \\ \hline
        insieme delle parti & $\subseteq$ & $\cup$ & $\cap$ & $\varnothing$ & $s$ & $s \setminus x$ \\
        valori di verità $\{\falso, \vero\}$ & $\falso$ prima di $\vero$ & OR & AND & $\falso$ & $\vero$ & NOT \\
    \end{tabular}
    \end{center}
    Per i valori di verità le proprietà 1--4 sono tautologie già viste nel Capitolo~1: commutatività, associatività, idempotenza e distributività
    di AND e OR, terzo escluso ($p \vee \neg p$) e non contraddizione ($\neg(p \wedge \neg p)$). Da qui il nome: George Boole studiò proprio
    l'«algebra della logica».
\end{esempio}

\begin{concetto}
\begin{Teorema}{Reticoli booleani e algebre di Boole}{reticoli-algebre-boole}
    \begin{prerequisiti}
        \dip{Teorema}{teo:proprieta-reticolari}
        \dip{Proposizione}{prop:neutri-reticolo}
        \dip{Lemma}{lem:ordine-da-operazioni}
        \dip{Teorema}{teo:reticoli-equivalenza}
        \dip{Teorema}{teo:complemento-unico}
    \end{prerequisiti}
    Se $(s, \rho)$ è un reticolo booleano, $(s, \vee_\rho, \wedge_\rho, \min(s), \max(s), {}')$ è un'algebra di Boole. Viceversa, se $(s, \vee, \wedge, 0, 1, {}')$ è
    un'algebra di Boole, la relazione $x \rho y \defsse x = x \wedge y$ rende $(s, \rho)$ un reticolo booleano con $\sup = \vee$, $\inf = \wedge$, $\min(s) = 0$,
    $\max(s) = 1$ e complemento ${}'$. I due passaggi sono uno l'inverso dell'altro.
\end{Teorema}
\begin{dimostrazione}
    \begin{caso}{Dal reticolo all'algebra}
        Le operazioni reticolari sono commutative, associative, idempotenti e soddisfano l'assorbimento (Teorema~\ref{teo:proprieta-reticolari}); $\max(s)$ è neutro
        per $\wedge$ e $\min(s)$ per $\vee$ (Proposizione~\ref{prop:neutri-reticolo}). La distributività è l'ipotesi, e $x'$ è un complemento di $x$: valgono 1--4.
    \end{caso}
    \begin{caso}{Dall'algebra al reticolo}
        Per i punti 1 e 2, $\vee$ e $\wedge$ sono commutative, associative e con l'assorbimento: per il Lemma~\ref{lem:ordine-da-operazioni} $(s, \rho)$ è un reticolo
        con $\inf\{x, y\} = x \wedge y$ e $\sup\{x, y\} = x \vee y$. Poiché $1$ è neutro per $\wedge$ e $0$ per $\vee$, $1 = \max(s)$ e $0 = \min(s)$ (Proposizione~\ref{prop:neutri-reticolo}):
        il reticolo è limitato. Per il punto 4 $x'$ è un complemento di $x$, e per il punto 3 il reticolo è distributivo.
    \end{caso}
    I passaggi tra ordine e operazioni sono inversi per il Teorema~\ref{teo:reticoli-equivalenza}; inoltre $0$, $1$ e ${}'$ sono determinati dall'ordine (sono il
    minimo, il massimo e l'unico complemento, Teorema~\ref{teo:complemento-unico}).
\end{dimostrazione}
\end{concetto}

\begin{spiegazione}
    Un'algebra di Boole è un reticolo booleano scritto «con le operazioni» invece che «con l'ordine»: le due descrizioni contengono le stesse
    informazioni. Ad esempio, i divisori di $30$ con la divisibilità diventano l'algebra di Boole $(D, \mathrm{mcm}, \mathrm{MCD}, 1, 30, x \mapsto 30/x)$, dove $D$ è
    l'insieme dei divisori:
    \begin{passi}
        \item lo «$0$» dell'algebra è il minimo, cioè il numero $1$, e l'«$1$» è il massimo, cioè il numero $30$;
        \item $x \vee x' = \mathrm{mcm}(x, 30/x) = 30$ e $x \wedge x' = \mathrm{MCD}(x, 30/x) = 1$, perché $x$ e $30/x$ non hanno fattori primi in comune;
        \item viceversa, dalla sola tavola di $\mathrm{MCD}$ si ritrova l'ordine: $x \mid y$ se e solo se $\mathrm{MCD}(x, y) = x$.
    \end{passi}
\end{spiegazione}

\begin{concetto}
\begin{Proposizione}{Leggi di De Morgan nei reticoli booleani}{demorgan-boole}
    \begin{prerequisiti}
        \dip{Definizione}{def:reticolo-booleano}
        \dip{Teorema}{teo:complemento-unico}
        \dip{Proposizione}{prop:ordine-operazioni}
        \dip{Teorema}{teo:dualita-reticoli}
    \end{prerequisiti}
    In un reticolo booleano, per ogni $x, y$: $\ (x')' = x$, $\ 0' = 1$, $\ 1' = 0$, $\ (x \wedge y)' = x' \vee y'\ $ e $\ (x \vee y)' = x' \wedge y'$.
\end{Proposizione}
\begin{dimostrazione}
    Ricordiamo che $0 \wedge z = 0$ e $1 \vee z = 1$ per ogni $z$ (Proposizione~\ref{prop:ordine-operazioni}). $x$ è un complemento di $x'$, $1$ di $0$ e $0$ di $1$: per l'unicità
    del complemento $(x')' = x$, $0' = 1$ e $1' = 0$. Mostriamo che $x' \vee y'$ è un complemento di $x \wedge y$:
    \[ (x \wedge y) \wedge (x' \vee y') \overset{\text{(D1)}}{=} (x \wedge x' \wedge y) \vee (x \wedge y \wedge y') = (0 \wedge y) \vee (x \wedge 0) = 0 \vee 0 = 0, \]
    \[ (x \wedge y) \vee (x' \vee y') \overset{\text{(D2)}}{=} (x \vee x' \vee y') \wedge (y \vee x' \vee y') = (1 \vee y') \wedge (1 \vee x') = 1 \wedge 1 = 1. \]
    Per l'unicità $(x \wedge y)' = x' \vee y'$. L'ultima uguaglianza è la duale.
\end{dimostrazione}
\end{concetto}

\begin{spiegazione}
    Sono le stesse leggi di De Morgan della logica (Proposizione~\ref{prop:demorgan}) e degli insiemi (Proposizione~\ref{prop:demorgan-ins}). La dimostrazione usa
    un'idea che torna spesso: per sapere chi è il complemento di qualcosa basta \textbf{indovinarlo e verificarlo}, perché il complemento è unico.
    Nei divisori di $30$ (dove il minimo è il numero $1$, il massimo è $30$ e $x' = 30/x$):
    \begin{passi}
        \item $(6 \wedge 10)' = 2' = 15$;
        \item $6' \vee 10' = 5 \vee 3 = 15$: le due cose coincidono.
    \end{passi}
\end{spiegazione}

% ──────────────────────────────────────────────────────────────
\section{Dagli anelli booleani ai reticoli booleani}\label{sec:anelli-reticoli}

Ricordiamo (Teorema~\ref{teo:booleani}) che un anello booleano è commutativo e che in esso $x + x = 0$, cioè $-x = x$, per ogni $x$.

\info{Una convenzione}{
    Come nel programma del corso, da qui in poi consideriamo solo anelli booleani, algebre di Boole e reticoli booleani con
    \textbf{almeno due elementi}, cioè con $0 \neq 1$. Ammettendo anche la struttura con un solo elemento ($0 = 1$) tutti i risultati restano veri,
    con l'unica differenza che negli enunciati sull'ordine $2^n$ si avrà $n \geq 0$ invece di $n \geq 1$: è solo una questione di coerenza nella notazione.
}

\begin{concetto}
\begin{Teorema}{L'ordine di un anello booleano}{anello-reticolo}
    \begin{prerequisiti}
        \dip{Definizione}{def:booleano}
        \dip{Teorema}{teo:booleani}
        \dip{Definizione}{def:tipi-relazioni}
        \dip{Definizione}{def:reticolo}
    \end{prerequisiti}
    Sia $(a, +, \cdot)$ un anello booleano e definiamo su $a$ la relazione $x \rho y \defsse xy = x$. Allora $\rho$ è un ordine, $(a, \rho)$ è un reticolo e
    \[ x \wedge y = xy \qquad\qquad x \vee y = x + y + xy. \]
\end{Teorema}
\begin{dimostrazione}
    Usiamo che $a$ è commutativo e che $z^2 = z$ e $z + z = 0$ per ogni $z$.
    \begin{caso}{$\rho$ è un ordine}
        Riflessiva: $xx = x$. Antisimmetrica: da $xy = x$ e $yx = y$ segue $x = xy = yx = y$. Transitiva: da $xy = x$ e $yz = y$ segue $xz = (xy)z = x(yz) = xy = x$.
    \end{caso}
    \begin{caso}{$xy$ è l'estremo inferiore}
        $(xy)x = x^2y = xy$ e $(xy)y = xy^2 = xy$: $xy$ è un minorante di $\{x, y\}$. Se $z$ è un minorante ($zx = z$ e $zy = z$), allora $z(xy) = (zx)y = zy = z$,
        cioè $z \rho xy$.
    \end{caso}
    \begin{caso}{$m = x + y + xy$ è l'estremo superiore}
        $xm = x^2 + xy + x^2y = x + (xy + xy) = x$ e $ym = xy + y^2 + xy^2 = y + (xy + xy) = y$: $m$ è un maggiorante di $\{x, y\}$. Se $z$ è un maggiorante
        ($xz = x$ e $yz = y$), allora $mz = xz + yz + x(yz) = x + y + xy = m$, cioè $m \rho z$.
    \end{caso}
\end{dimostrazione}
\end{concetto}

\begin{spiegazione}
    Nell'anello $(\mcP(s), \Delta, \cap)$: $x \cap y = x$ vuol dire $x \subseteq y$, quindi l'ordine è l'inclusione; l'inf è $x \cap y$ e il sup è
    $x \,\Delta\, y \,\Delta\, (x \cap y)$, che è proprio $x \cup y$:
    \begin{center}
    \begin{tikzpicture}
        \begin{scope}
            \fill[primary!25, even odd rule] (0,0) circle (0.75) (0.9,0) circle (0.75);
            \draw[black!60] (0,0) circle (0.75) (0.9,0) circle (0.75);
            \node[font=\small] at (0.45,-1.1) {$x \,\Delta\, y$};
        \end{scope}
        \node at (2.3,0) {$\Delta$};
        \begin{scope}[xshift=3.6cm]
            \begin{scope} \clip (0,0) circle (0.75); \fill[primary!25] (0.9,0) circle (0.75); \end{scope}
            \draw[black!60] (0,0) circle (0.75) (0.9,0) circle (0.75);
            \node[font=\small] at (0.45,-1.1) {$x \cap y$};
        \end{scope}
        \node at (5.9,0) {$=$};
        \begin{scope}[xshift=7.2cm]
            \fill[primary!25] (0,0) circle (0.75) (0.9,0) circle (0.75);
            \draw[black!60] (0,0) circle (0.75) (0.9,0) circle (0.75);
            \node[font=\small] at (0.45,-1.1) {$x \cup y$};
        \end{scope}
    \end{tikzpicture}
    \end{center}
    La differenza simmetrica «dimentica» la parte comune; aggiungendo con $\Delta$ la parte comune la si ritrova, e si ottiene l'unione.
\end{spiegazione}

\begin{concetto}
\begin{Proposizione}{Ogni anello booleano è un reticolo booleano}{anello-reticolo-booleano}
    \begin{prerequisiti}
        \dip{Teorema}{teo:anello-reticolo}
        \dip{Definizione}{def:reticolo-booleano}
        \dip{Definizione}{def:distributivo}
    \end{prerequisiti}
    Il reticolo $(a, \rho)$ del Teorema~\ref{teo:anello-reticolo} è booleano, con minimo $0$, massimo $1$ e complemento $x' = 1 + x$.
\end{Proposizione}
\begin{dimostrazione}
    \begin{caso}{Limitato e complementato}
        $0x = 0$ e $x1 = x$, cioè $0 \rho x \rho 1$: $0$ e $1$ sono il minimo e il massimo. Inoltre $x \wedge (1 + x) = x + x^2 = x + x = 0$ e
        $x \vee (1 + x) = x + (1 + x) + x(1 + x) = (x + x) + 1 + (x + x) = 1$: $1 + x$ è un complemento di $x$.
    \end{caso}
    \begin{caso}{Distributivo}
        (D1): $x \wedge (y \vee z) = x(y + z + yz) = xy + xz + xyz$ e $(x \wedge y) \vee (x \wedge z) = xy + xz + (xy)(xz) = xy + xz + xyz$. (D2): sviluppando il prodotto
        $(x \vee y) \wedge (x \vee z) = (x + y + xy)(x + z + xz)$ e semplificando le coppie di termini uguali ($w + w = 0$) si ottiene $x + yz + xyz = x \vee (y \wedge z)$.
    \end{caso}
\end{dimostrazione}
\end{concetto}

\begin{spiegazione}
    Il complemento è «sommare $1$»: in $(\mcP(s), \Delta, \cap)$ è $s \,\Delta\, x = s \setminus x$, e sulle stringhe di bit (che vedremo alla fine del
    capitolo) è la negazione di ogni bit. Nell'anello più piccolo, $\mbZ_2 = \{0, 1\}$, l'ordine è $0 \rho 1$ e il complemento scambia $0$ e $1$: è il
    NOT dei valori di verità.
\end{spiegazione}

% ──────────────────────────────────────────────────────────────
\newpage
\section{Dai reticoli booleani agli anelli booleani}\label{sec:reticoli-anelli}

Ora il viceversa: in un reticolo booleano definiamo una somma che imita la differenza simmetrica.

\begin{concetto}
\begin{Lemma}{La somma booleana}{somma-booleana}
    \begin{prerequisiti}
        \dip{Definizione}{def:reticolo-booleano}
        \dip{Proposizione}{prop:demorgan-boole}
        \dip{Proposizione}{prop:ordine-operazioni}
        \dip{Teorema}{teo:proprieta-reticolari}
    \end{prerequisiti}
    Sia $(s, \rho)$ un reticolo booleano e poniamo $x + y := (x \wedge y') \vee (x' \wedge y)$. Allora, per ogni $x, y, z \in s$:
    \begin{enumerate}
        \item $x + y = y + x$, \ $x + 0 = x$ \ e \ $x + x = 0$;
        \item $(x + y)' = (x \wedge y) \vee (x' \wedge y')$;
        \item $(x + y) + z = (x \wedge y' \wedge z') \vee (x' \wedge y \wedge z') \vee (x' \wedge y' \wedge z) \vee (x \wedge y \wedge z)$, e quindi $+$ è associativa.
    \end{enumerate}
\end{Lemma}
\begin{dimostrazione}
    Useremo che $0 \wedge z = 0$, $1 \wedge z = z$ e $0 \vee z = z$ (Proposizione~\ref{prop:ordine-operazioni}).
    \begin{caso}{Punto 1}
        Lo scambio di $x$ e $y$ scambia i due termini di $\vee$. $x + 0 = (x \wedge 0') \vee (x' \wedge 0) = (x \wedge 1) \vee 0 = x$ e $x + x = (x \wedge x') \vee (x' \wedge x) = 0 \vee 0 = 0$.
    \end{caso}
    \begin{caso}{Punto 2}
        Per De Morgan $(x + y)' = (x \wedge y')' \wedge (x' \wedge y)' = (x' \vee y) \wedge (x \vee y')$; distribuendo,
        $= (x' \wedge x) \vee (x' \wedge y') \vee (y \wedge x) \vee (y \wedge y') = (x \wedge y) \vee (x' \wedge y')$.
    \end{caso}
    \begin{caso}{Punto 3}
        $(x + y) + z = \big((x + y) \wedge z'\big) \vee \big((x + y)' \wedge z\big)$; sostituendo la definizione e il punto 2 e distribuendo si ottengono i quattro termini.
        L'espressione non cambia permutando $x$, $y$, $z$ (i quattro termini si scambiano tra loro), quindi $(x + y) + z = (y + z) + x = x + (y + z)$.
    \end{caso}
\end{dimostrazione}
\end{concetto}

\begin{spiegazione}
    I quattro termini del punto 3 descrivono gli elementi che stanno in \textbf{esattamente uno} dei tre insiemi oppure in \textbf{tutti e tre}: cioè in un
    numero \textbf{dispari} di essi. Questa descrizione non dipende dall'ordine dei tre insiemi, ed ecco l'associatività. Nel diagramma di Venn
    $x \,\Delta\, y \,\Delta\, z$ è la parte colorata:
    \begin{center}
    \begin{tikzpicture}
        \fill[primary!25, even odd rule] (0,0) circle (0.85) (1.1,0) circle (0.85) (0.55,-0.95) circle (0.85);
        \draw[black!60] (0,0) circle (0.85) (1.1,0) circle (0.85) (0.55,-0.95) circle (0.85);
        \node at (-0.4,0.35) {$x$}; \node at (1.5,0.35) {$y$}; \node at (0.55,-1.45) {$z$};
        \node[font=\small, align=left, anchor=west] at (2.6,-0.45) {colorati: i punti che stanno in $1$ o in $3$ dei tre insiemi;\\
            bianchi: quelli che stanno in $0$ o in $2$.\\[1mm]
            È come lo \texttt{XOR} di tre bit: vale $1$ se e solo se un numero dispari di bit vale $1$.};
    \end{tikzpicture}
    \end{center}
\end{spiegazione}

\begin{concetto}
\begin{Teorema}{Ogni reticolo booleano è un anello booleano}{reticolo-anello}
    \begin{prerequisiti}
        \dip{Lemma}{lem:somma-booleana}
        \dip{Proposizione}{prop:neutri-reticolo}
        \dip{Definizione}{def:anello}
        \dip{Definizione}{def:booleano}
    \end{prerequisiti}
    Sia $(s, \rho)$ un reticolo booleano. Con $x + y := (x \wedge y') \vee (x' \wedge y)$ e $x \cdot y := x \wedge y$, $(s, +, \cdot)$ è un anello booleano, con zero
    $0 = \min(s)$, unità $1 = \max(s)$ e $-x = x$ per ogni $x$.
\end{Teorema}
\begin{dimostrazione}
    \begin{caso}{$(s, +)$ è un gruppo abeliano e $(s, \cdot)$ un monoide commutativo}
        Per il Lemma~\ref{lem:somma-booleana} $+$ è associativa e commutativa, $0$ è neutro e ogni $x$ è l'opposto di se stesso. $\wedge$ è associativa, commutativa
        e ha neutro $1 = \max(s)$ (Proposizione~\ref{prop:neutri-reticolo}); inoltre $x \cdot x = x \wedge x = x$.
    \end{caso}
    \begin{caso}{Distributività}
        Per la distributività del reticolo, $x(y + z) = x \wedge \big((y \wedge z') \vee (y' \wedge z)\big) = (x \wedge y \wedge z') \vee (x \wedge y' \wedge z)$. D'altra parte, per De Morgan,
        \[ xy + xz = \big((x \wedge y) \wedge (x' \vee z')\big) \vee \big((x' \vee y') \wedge (x \wedge z)\big) = (x \wedge y \wedge z') \vee (x \wedge y' \wedge z), \]
        perché i termini con $x \wedge x'$ valgono $0$. La distributività a destra segue dalla commutatività di $\cdot$.
    \end{caso}
    Quindi $(s, +, \cdot)$ è un anello commutativo unitario in cui $x^2 = x$: un anello booleano.
\end{dimostrazione}
\end{concetto}

\begin{spiegazione}
    Calcoliamo nei divisori di $30$, dove $x \cdot y = \mathrm{MCD}(x, y)$ e $x' = 30/x$:
    \begin{passi}
        \item $6 + 10 = (6 \wedge 10') \vee (6' \wedge 10) = (6 \wedge 3) \vee (5 \wedge 10) = 3 \vee 5 = 15$;
        \item $6 \cdot 10 = \mathrm{MCD}(6, 10) = 2$;
        \item controllo con i fattori primi: $6 \leftrightarrow \{2, 3\}$ e $10 \leftrightarrow \{2, 5\}$; $\{2, 3\} \,\Delta\, \{2, 5\} = \{3, 5\} \leftrightarrow 15$ e $\{2, 3\} \cap \{2, 5\} = \{2\} \leftrightarrow 2$.
    \end{passi}
    La somma tiene i fattori primi che stanno in \textbf{uno solo} dei due numeri, il prodotto quelli che stanno in \textbf{entrambi}.
\end{spiegazione}

\begin{concetto}
\begin{Teorema}{Anelli booleani e reticoli booleani si corrispondono}{boole-inverse}
    \begin{prerequisiti}
        \dip{Teorema}{teo:anello-reticolo}
        \dip{Proposizione}{prop:anello-reticolo-booleano}
        \dip{Teorema}{teo:reticolo-anello}
        \dip{Proposizione}{prop:ordine-operazioni}
    \end{prerequisiti}
    Le due costruzioni (dall'anello booleano al reticolo booleano e viceversa) sono una l'inversa dell'altra: partendo da un anello booleano, passando
    al reticolo e tornando all'anello si ritrovano le stesse operazioni; partendo da un reticolo booleano e facendo andata e ritorno si ritrova lo stesso ordine.
\end{Teorema}
\begin{dimostrazione}
    \begin{caso}{Anello $\to$ reticolo $\to$ anello}
        Nel reticolo dell'anello $(a, +, \cdot)$ è $x \wedge y = xy$: il nuovo prodotto è quello vecchio. La nuova somma è $(x \wedge y') \vee (x' \wedge y) = u \vee v$ con
        $u = x(1 + y) = x + xy$ e $v = (1 + x)y = y + xy$. Ora $uv = xy + xy + xy + xy = 0$, quindi $u \vee v = u + v + uv = x + y + (xy + xy) = x + y$.
    \end{caso}
    \begin{caso}{Reticolo $\to$ anello $\to$ reticolo}
        Nell'anello del reticolo $(s, \rho)$ è $xy = x \wedge y$; il nuovo ordine è $xy = x$, cioè $x \wedge y = x$, che equivale a $x \rho y$ (Proposizione~\ref{prop:ordine-operazioni}).
    \end{caso}
\end{dimostrazione}
\end{concetto}

\begin{spiegazione}
    Ecco il quadro completo. Ogni freccia è una costruzione, e andata e ritorno riportano sempre al punto di partenza:
    \begin{center}
    \begin{tikzpicture}[scatola/.style={draw=primary!70!black, fill=primary!6, rounded corners=3pt, inner sep=5pt, align=center, font=\small}]
        \node[scatola] (R) at (0,0) {\textbf{reticolo booleano}\\ $(s, \rho)$};
        \node[scatola] (B) at (-5,-2.3) {\textbf{algebra di Boole}\\ $(s, \vee, \wedge, 0, 1, {}')$};
        \node[scatola] (A) at (5,-2.3) {\textbf{anello booleano}\\ $(s, +, \cdot)$};
        \draw[-{Stealth}, line width=0.9pt, primary, transform canvas={shift={(-0.07,0.13)}}] (R) -- node[above left=-0.5mm] {\numcerchio{1}} (B);
        \draw[-{Stealth}, line width=0.9pt, primary, transform canvas={shift={(0.07,-0.13)}}] (B) -- node[below right=-0.5mm] {\numcerchio{2}} (R);
        \draw[-{Stealth}, line width=0.9pt, primary, transform canvas={shift={(0.07,0.13)}}] (R) -- node[above right=-0.5mm] {\numcerchio{3}} (A);
        \draw[-{Stealth}, line width=0.9pt, primary, transform canvas={shift={(-0.07,-0.13)}}] (A) -- node[below left=-0.5mm] {\numcerchio{4}} (R);
    \end{tikzpicture}
    \end{center}
    \begin{center}
    \renewcommand{\arraystretch}{1.3}
    \begin{tabular}{c|l|l}
        & \textbf{passaggio} & \textbf{formule} \\ \hline
        \numcerchio{1} & dall'ordine alle operazioni & $\vee = \sup$, \ $\wedge = \inf$, \ $0 = \min$, \ $1 = \max$, \ ${}' = $ complemento \\
        \numcerchio{2} & dalle operazioni all'ordine & $x \rho y \leftrightarrow x = x \wedge y$ \\
        \numcerchio{3} & dal reticolo all'anello & $xy := x \wedge y$, \ $x + y := (x \wedge y') \vee (x' \wedge y)$ \\
        \numcerchio{4} & dall'anello al reticolo & $x \rho y \leftrightarrow xy = x$; \ allora $x \wedge y = xy$, \ $x \vee y = x + y + xy$, \ $x' = 1 + x$ \\
    \end{tabular}
    \end{center}
    Nell'insieme delle parti: $(\mcP(s), \subseteq)$, $(\mcP(s), \cup, \cap, \varnothing, s, s \setminus \cdot)$ e $(\mcP(s), \Delta, \cap)$ sono lo stesso oggetto. Infatti
    $(x \cap (s \setminus y)) \cup ((s \setminus x) \cap y) = (x \setminus y) \cup (y \setminus x) = x \,\Delta\, y$.
\end{spiegazione}

\begin{concetto}
\begin{Proposizione}{Isomorfismi di anelli booleani e di reticoli booleani}{isomorfismi-booleani}
    \begin{prerequisiti}
        \dip{Teorema}{teo:boole-inverse}
        \dip{Teorema}{teo:isomorfismi-reticoli}
        \dip{Teorema}{teo:complemento-unico}
        \dip{Definizione}{def:omomorfismo-anelli}
    \end{prerequisiti}
    Siano $(a, +, \cdot)$ e $(a', +, \cdot)$ anelli booleani, con i reticoli associati $(a, \rho)$ e $(a', \rho')$, e sia $f \colon a \to a'$ biettiva. Allora $f$ è un isomorfismo di
    anelli se e solo se è un isomorfismo di insiemi ordinati tra $(a, \rho)$ e $(a', \rho')$.
\end{Proposizione}
\begin{dimostrazione}
    \begin{caso}{$(\rightarrow)$}
        $x \rho y \leftrightarrow xy = x \leftrightarrow f(xy) = f(x) \leftrightarrow f(x)f(y) = f(x) \leftrightarrow f(x) \rho' f(y)$, dove il secondo $\leftrightarrow$ usa l'iniettività.
    \end{caso}
    \begin{caso}{$(\leftarrow)$}
        $f$ conserva $\wedge$ e $\vee$ (Teorema~\ref{teo:isomorfismi-reticoli}) e manda il minimo nel minimo e il massimo nel massimo. Quindi da $x \wedge x' = 0$ e $x \vee x' = 1$
        segue che $f(x')$ è il complemento di $f(x)$: $f(x') = f(x)'$. Per il Teorema~\ref{teo:boole-inverse} le operazioni dell'anello si scrivono con $\wedge$, $\vee$ e ${}'$:
        $f(xy) = f(x \wedge y) = f(x) \wedge f(y) = f(x)f(y)$ e $f(x + y) = f\big((x \wedge y') \vee (x' \wedge y)\big) = (f(x) \wedge f(y)') \vee (f(x)' \wedge f(y)) = f(x) + f(y)$.
    \end{caso}
\end{dimostrazione}
\end{concetto}

\begin{spiegazione}
    Le operazioni di un anello booleano si scrivono con l'ordine, e l'ordine si scrive con le operazioni: quindi chi conserva le une conserva
    anche l'altro, e viceversa. In pratica, per vedere se due anelli booleani finiti sono isomorfi basta confrontare i diagrammi di Hasse dei
    reticoli associati (Teorema~\ref{teo:hasse-isomorfismo}).
\end{spiegazione}

\begin{esempio}{I divisori di 210}
    Sia $r = \{n \in \mbN \mid n \mid 210\}$, con $210 = 2 \cdot 3 \cdot 5 \cdot 7$. È un sottoreticolo di $(\mbN, \mid)$, quindi distributivo, e $x' = 210/x$ è un complemento: è un
    reticolo booleano con $16$ elementi, minimo $1$ e massimo $210$.
    \begin{center}
    \begin{tikzpicture}[hasse, x=1cm, y=1cm]
        \node[hv] (1) at (0,0) {$1$};
        \node[hv evid] (2) at (-2.4,1.15) {$2$}; \node[hv] (3) at (-0.8,1.15) {$3$}; \node[hv] (5) at (0.8,1.15) {$5$}; \node[hv evid] (7) at (2.4,1.15) {$7$};
        \node[hv evid] (6) at (-3.75,2.3) {$6$}; \node[hv evid] (10) at (-2.25,2.3) {$10$}; \node[hv evid2] (14) at (-0.75,2.3) {$14$};
        \node[hv evid2] (15) at (0.75,2.3) {$15$}; \node[hv] (21) at (2.25,2.3) {$21$}; \node[hv] (35) at (3.75,2.3) {$35$};
        \node[hv evid2] (30) at (-2.4,3.45) {$30$}; \node[hv] (42) at (-0.8,3.45) {$42$}; \node[hv] (70) at (0.8,3.45) {$70$}; \node[hv] (105) at (2.4,3.45) {$105$};
        \node[hv] (210) at (0,4.6) {$210$};
        \foreach \a/\b in {1/2,1/3,1/5,1/7, 2/6,2/10,2/14, 3/6,3/15,3/21, 5/10,5/15,5/35, 7/14,7/21,7/35,
                           6/30,6/42, 10/30,10/70, 14/42,14/70, 15/30,15/105, 21/42,21/105, 35/70,35/105,
                           30/210,42/210,70/210,105/210}
            \draw[hlato] (\a) -- (\b);
        \node[font=\small, align=left, anchor=west] at (4.6,2.3) {
            $2 \wedge 6 = \mathrm{MCD}(2, 6) = 2$\\[0.6mm]
            $2 \vee 7 = \mathrm{mcm}(2, 7) = 14$\\[0.6mm]
            $7' = 210 / 7 = 30$\\[0.6mm]
            $6 + 10 = (6 \wedge 35) \vee (35 \wedge 10) = 3 \vee 5 = 15$\\[0.6mm]
            $6 \cdot 10 = 6 \wedge 10 = 2$\\[1.5mm]
            in arancione i dati, in verde i risultati};
    \end{tikzpicture}
    \end{center}
    L'applicazione che manda ogni divisore nell'insieme dei suoi fattori primi è un isomorfismo tra $(r, \mid)$ e $(\mcP(\{2, 3, 5, 7\}), \subseteq)$: il diagramma è un
    «ipercubo» a quattro dimensioni. Invece i divisori di $12$ non formano un reticolo booleano: $2$ non ha complemento (Esempio~\ref{es:divisori-30-12}).
\end{esempio}

% ──────────────────────────────────────────────────────────────
\section{Il teorema di Stone}\label{sec:stone}

L'insieme delle parti non è solo un esempio di anello booleano: in un senso preciso è l'\textbf{unico} esempio.

\begin{concetto}
\begin{Teorema}{Teorema di rappresentazione di Stone}{stone}
    \begin{prerequisiti}
        \dip{Definizione}{def:booleano}
        \dip{Definizione}{def:omomorfismo-anelli}
    \end{prerequisiti}
    \senzadim
    \begin{enumerate}
        \item (Caso generale) Ogni anello booleano è isomorfo a un sottoanello di $(\mcP(s), \Delta, \cap)$ per un opportuno insieme $s$.
        \item (Caso finito) Ogni anello booleano finito è isomorfo a $(\mcP(s), \Delta, \cap)$ per un opportuno insieme finito $s$.
    \end{enumerate}
\end{Teorema}
\end{concetto}

\begin{spiegazione}
    Non vediamo la dimostrazione, ma l'idea del caso finito si vede sul diagramma. Gli \textbf{atomi} sono gli elementi che coprono il minimo;
    ogni elemento è il sup degli atomi che stanno sotto di lui, e l'isomorfismo manda ogni elemento nell'insieme dei suoi atomi:
    \begin{center}
    \begin{tikzpicture}[hasse]
        \node[hv] (1) at (0,0) {$1$};
        \node[hv evid] (2) at (-1.2,1.1) {$2$}; \node[hv evid] (3) at (0,1.1) {$3$}; \node[hv evid] (5) at (1.2,1.1) {$5$};
        \node[hv] (6) at (-1.2,2.2) {$6$}; \node[hv] (10) at (0,2.2) {$10$}; \node[hv] (15) at (1.2,2.2) {$15$};
        \node[hv] (30) at (0,3.3) {$30$};
        \draw[hlato] (1) -- (2); \draw[hlato] (1) -- (3); \draw[hlato] (1) -- (5);
        \draw[hlato] (2) -- (6); \draw[hlato] (2) -- (10); \draw[hlato] (3) -- (6); \draw[hlato] (3) -- (15); \draw[hlato] (5) -- (10); \draw[hlato] (5) -- (15);
        \draw[hlato] (6) -- (30); \draw[hlato] (10) -- (30); \draw[hlato] (15) -- (30);
        \node[font=\small, align=left, anchor=west] at (2.2,1.65) {nei divisori di $30$ gli atomi (in arancione) sono $2$, $3$, $5$;\\
            $6 = 2 \vee 3 \mapsto \{2, 3\}$, \ $30 = 2 \vee 3 \vee 5 \mapsto \{2, 3, 5\}$, \ $1 \mapsto \varnothing$:\\
            così i divisori di $30$ diventano le parti di $s = \{2, 3, 5\}$.};
    \end{tikzpicture}
    \end{center}
    Nel caso infinito gli atomi possono non bastare (ci sono anelli booleani infiniti senza atomi), ed è per questo che si ottiene solo un
    sottoanello di un insieme delle parti.
\end{spiegazione}

\begin{concetto}
\begin{Corollario}{Teorema di Stone per i reticoli booleani}{stone-reticoli}
    \begin{prerequisiti}
        \dip{Teorema}{teo:stone}
        \dip{Teorema}{teo:boole-inverse}
        \dip{Proposizione}{prop:isomorfismi-booleani}
        \dip{Proposizione}{prop:caratt-inclusione}
    \end{prerequisiti}
    Ogni reticolo booleano finito $(s, \rho)$ è isomorfo, come insieme ordinato, a $(\mcP(t), \subseteq)$ per un opportuno insieme finito $t$.
\end{Corollario}
\begin{dimostrazione}
    L'anello booleano associato a $(s, \rho)$ è isomorfo a $(\mcP(t), \Delta, \cap)$ per un insieme finito $t$ (Teorema~\ref{teo:stone}). Per la Proposizione~\ref{prop:isomorfismi-booleani}
    lo stesso isomorfismo è un isomorfismo tra i reticoli associati: quello di $s$ è $(s, \rho)$ (Teorema~\ref{teo:boole-inverse}) e quello di $\mcP(t)$ ha l'ordine
    $x \cap y = x$, cioè $x \subseteq y$ (Proposizione~\ref{prop:caratt-inclusione}).
\end{dimostrazione}
\end{concetto}

\begin{spiegazione}
    Ogni reticolo booleano finito «è» un cubo: un segmento ($2$ elementi), un quadrato ($4$), un cubo ($8$), un ipercubo ($16$, come i divisori
    di $210$), e così via. Ad esempio $(\mbZ_2, \rho)$, con $0 \rho 1$, è il segmento $\mcP(\{1\})$, e i divisori di $6$ sono il quadrato $\mcP(\{2, 3\})$.
\end{spiegazione}

\begin{concetto}
\begin{Lemma}{Insiemi equipotenti hanno anelli delle parti isomorfi}{biezione-parti}
    \begin{prerequisiti}
        \dip{Definizione}{def:immagine-antiimmagine}
        \dip{Proposizione}{prop:isomorfismi-booleani}
        \dip{Teorema}{teo:biettiva-inversa}
    \end{prerequisiti}
    Se $g \colon s \to s'$ è biettiva, allora $G \colon x \in \mcP(s) \mapsto \overrightarrow{g}(x) \in \mcP(s')$ è un isomorfismo tra $(\mcP(s), \Delta, \cap)$ e $(\mcP(s'), \Delta, \cap)$.
\end{Lemma}
\begin{dimostrazione}
    $G$ è biettiva: la sua inversa è $y \mapsto \overrightarrow{g^{-1}}(y)$. Se $x \subseteq y$, ovviamente $G(x) \subseteq G(y)$. Viceversa, se $G(x) \subseteq G(y)$ e $z \in x$, allora $g(z) = g(w)$ per
    un $w \in y$, e per l'iniettività $z = w \in y$: quindi $x \subseteq y$. Dunque $G$ è un isomorfismo tra $(\mcP(s), \subseteq)$ e $(\mcP(s'), \subseteq)$, cioè tra i reticoli
    associati ai due anelli, e per la Proposizione~\ref{prop:isomorfismi-booleani} è un isomorfismo di anelli.
\end{dimostrazione}
\end{concetto}

\begin{spiegazione}
    Rinominare gli elementi di $s$ non cambia la struttura di $\mcP(s)$: conta solo quanti sono. Con $g(a) = 1$, $g(b) = 2$, la parte $\{a, b\}$ diventa
    $\{1, 2\}$, $\{a\}$ diventa $\{1\}$, e $\{a\} \,\Delta\, \{a, b\} = \{b\}$ diventa $\{1\} \,\Delta\, \{1, 2\} = \{2\}$.
\end{spiegazione}

\begin{concetto}
\begin{Corollario}{Conseguenze del teorema di Stone}{conseguenze-stone}
    \begin{prerequisiti}
        \dip{Teorema}{teo:stone}
        \dip{Corollario}{cor:stone-reticoli}
        \dip{Lemma}{lem:biezione-parti}
        \dip{Corollario}{cor:cardinalita-parti}
        \dip{Teorema}{teo:inverso-isomorfismo}
    \end{prerequisiti}
    \begin{enumerate}
        \item Ogni anello booleano finito e ogni reticolo booleano finito hanno ordine $2^n$ con $n \geq 1$; per ogni $n \geq 1$ ce n'è uno di ordine $2^n$.
        \item Due anelli booleani finiti sono isomorfi se e solo se hanno lo stesso ordine (e lo stesso vale per i reticoli booleani finiti).
    \end{enumerate}
\end{Corollario}
\begin{dimostrazione}
    \begin{caso}{Punto 1}
        Un anello booleano finito $a$ è in biezione con $\mcP(s)$, $s$ finito (Teorema~\ref{teo:stone}), quindi $|a| = 2^{|s|}$ (Corollario~\ref{cor:cardinalita-parti}); $s \neq \varnothing$ perché $a$
        ha almeno due elementi. Per i reticoli si usa il Corollario~\ref{cor:stone-reticoli}. Viceversa $(\mcP(I_n), \Delta, \cap)$ ha ordine $2^n$.
    \end{caso}
    \begin{caso}{Punto 2}
        Un isomorfismo è una biezione, quindi due anelli isomorfi hanno lo stesso ordine. Viceversa, se $a \cong \mcP(s)$ e $a' \cong \mcP(s')$ con $2^{|s|} = |a| = |a'| = 2^{|s'|}$,
        allora $|s| = |s'|$ e c'è una biezione $g \colon s \to s'$. Per il Lemma~\ref{lem:biezione-parti} $\mcP(s) \cong \mcP(s')$, e componendo gli isomorfismi (e le loro
        inverse, Teorema~\ref{teo:inverso-isomorfismo}) si ottiene $a \cong a'$. Per i reticoli si ragiona allo stesso modo.
    \end{caso}
\end{dimostrazione}
\end{concetto}

\begin{spiegazione}
    Non esistono anelli booleani (né reticoli booleani) con $3$, $5$, $6$, $7$, $9$, \dots{} elementi, e quelli con lo stesso numero di elementi sono tutti
    «uguali a meno dei nomi». Ad esempio i divisori di $30$, $(\mcP(\{a, b, c\}), \subseteq)$ e i divisori di $105 = 3 \cdot 5 \cdot 7$ sono tre reticoli booleani di $8$ elementi:
    sono isomorfi. Un reticolo con $12$ elementi, come i divisori di $60$, non può essere booleano.
\end{spiegazione}

% ──────────────────────────────────────────────────────────────
\section{Funzioni caratteristiche e operazioni bit a bit}\label{sec:bit-a-bit}

Le funzioni caratteristiche (Definizione~\ref{def:funzione-caratteristica}) prendono valori in $\{0, 1\}$. Identifichiamo $\{0, 1\}$ con $\mbZ_2 = \{[0]_2, [1]_2\}$,
scrivendo semplicemente $0$ e $1$: le sue operazioni sono
\begin{center}
\renewcommand{\arraystretch}{1.15}
\begin{tabular}{c|cc}
    $+$ & $0$ & $1$ \\ \hline
    $0$ & $0$ & $1$ \\
    $1$ & $1$ & $0$ \\
\end{tabular}
\qquad\qquad
\begin{tabular}{c|cc}
    $\cdot$ & $0$ & $1$ \\ \hline
    $0$ & $0$ & $0$ \\
    $1$ & $0$ & $1$ \\
\end{tabular}
\end{center}
cioè la somma è lo \texttt{XOR} e il prodotto è l'\texttt{AND}. $\mbZ_2$ è un campo (Teorema~\ref{teo:zm-campo}), e in esso $0^2 = 0$ e $1^2 = 1$: è un anello booleano.

\dfn[z2n]{L'anello $\mbZ_2^n$ delle stringhe di bit}{
    Sia $n \geq 1$. Su $\mbZ_2^n := \mathrm{Map}(I_n, \mbZ_2)$ definiamo le operazioni \textbf{componente per componente}:
    \[ (f + g)(i) := f(i) + g(i), \qquad\qquad (f \cdot g)(i) := f(i) \cdot g(i) \qquad \text{per ogni } i \in I_n. \]
    Scriviamo $f$ come la \textbf{stringa di bit} $f(1)\,f(2) \dots f(n)$: ad esempio, con $n = 3$, la funzione con $f(1) = 0$, $f(2) = 1$, $f(3) = 1$ è la stringa \texttt{011}.
}

\begin{concetto}
\begin{Proposizione}{Le stringhe di bit formano un anello booleano}{z2n-booleano}
    \begin{prerequisiti}
        \dip{Definizione}{def:z2n}
        \dip{Teorema}{teo:zm-campo}
        \dip{Teorema}{teo:conta-map}
        \dip{Definizione}{def:booleano}
    \end{prerequisiti}
    $(\mbZ_2^n, +, \cdot)$ è un anello booleano con $2^n$ elementi; lo zero è la stringa \texttt{00\dots0} e l'unità è la stringa \texttt{11\dots1}.
\end{Proposizione}
\begin{dimostrazione}
    Ogni proprietà richiesta (associatività, commutatività, distributività, neutri, opposti, $f^2 = f$) è un'uguaglianza tra funzioni; calcolandola in un
    qualsiasi $i \in I_n$ diventa la stessa uguaglianza in $\mbZ_2$, che è vera perché $\mbZ_2$ è un anello booleano. Due funzioni che coincidono in ogni $i$ sono uguali,
    quindi le proprietà valgono in $\mbZ_2^n$. Lo zero e l'unità sono le funzioni costanti $0$ e $1$, e $|\mathrm{Map}(I_n, \mbZ_2)| = 2^n$ (Teorema~\ref{teo:conta-map}).
\end{dimostrazione}
\end{concetto}

\begin{spiegazione}
    Le operazioni si fanno «in colonna», bit per bit, senza riporti:
    \begin{center}
    \renewcommand{\arraystretch}{1.1}
    \begin{tabular}{r@{\ \ }l}
        & \texttt{0011} \\
        $+$ & \texttt{0101} \\ \hline
        & \texttt{0110} \\
    \end{tabular}
    \qquad\qquad\qquad
    \begin{tabular}{r@{\ \ }l}
        & \texttt{0011} \\
        $\cdot$ & \texttt{0101} \\ \hline
        & \texttt{0001} \\
    \end{tabular}
    \end{center}
    Ogni colonna è un calcolo in $\mbZ_2$: per questo valgono in $\mbZ_2^n$ tutte le proprietà di $\mbZ_2$.
\end{spiegazione}

\begin{concetto}
\begin{Teorema}{Parti di $I_n$ e stringhe di bit}{caratteristiche-z2n}
    \begin{prerequisiti}
        \dip{Teorema}{teo:parti-caratteristiche}
        \dip{Proposizione}{prop:appartenenza}
        \dip{Proposizione}{prop:z2n-booleano}
        \dip{Definizione}{def:omomorfismo-anelli}
    \end{prerequisiti}
    Sia $n \geq 1$. La funzione $\chi \colon t \in \mcP(I_n) \mapsto \chi_{t, I_n} \in \mbZ_2^n$ è un isomorfismo di anelli tra $(\mcP(I_n), \Delta, \cap)$ e $(\mbZ_2^n, +, \cdot)$.
\end{Teorema}
\begin{dimostrazione}
    $\chi$ è biettiva per il Teorema~\ref{teo:parti-caratteristiche}. Siano $a, b \subseteq I_n$ e $i \in I_n$. Per la Proposizione~\ref{prop:appartenenza}
    \[ \chi_{a \Delta b}(i) = 1 \leftrightarrow i \in a \,\Delta\, b \leftrightarrow (i \in a) \XOR (i \in b) \leftrightarrow \text{esattamente uno tra } \chi_a(i), \chi_b(i) \text{ vale } 1 \leftrightarrow \chi_a(i) + \chi_b(i) = 1. \]
    Poiché entrambi i membri valgono $0$ o $1$, $\chi_{a \Delta b}(i) = \chi_a(i) + \chi_b(i)$ per ogni $i$, cioè $\chi_{a \Delta b} = \chi_a + \chi_b$. Allo stesso modo
    $\chi_{a \cap b}(i) = 1 \leftrightarrow i \in a \text{ e } i \in b \leftrightarrow \chi_a(i) = \chi_b(i) = 1 \leftrightarrow \chi_a(i)\chi_b(i) = 1$, quindi $\chi_{a \cap b} = \chi_a \chi_b$.
\end{dimostrazione}
\end{concetto}

\begin{spiegazione}
    Una parte di $I_n$ si scrive come stringa di bit mettendo $1$ nelle posizioni dei suoi elementi. Con $n = 6$: $I_6 \mapsto \texttt{111111}$,
    $\{2, 3, 5\} \mapsto \texttt{011010}$, $\{1, 5, 6\} \mapsto \texttt{100011}$. Allora le operazioni sugli insiemi diventano le operazioni \textbf{bit a bit} dei linguaggi
    di programmazione (in C, Java, Python):
    \begin{center}
    \renewcommand{\arraystretch}{1.25}
    \begin{tabular}{l|l|l|l}
        insiemi & anello booleano & bit a bit & esempio con \texttt{011010} e \texttt{100011} \\ \hline
        $a \cap b$ & $ab$ & \texttt{a \& b} (\texttt{AND}) & \texttt{000010} $= \{5\}$ \\
        $a \,\Delta\, b$ & $a + b$ & \texttt{a \textasciicircum{} b} (\texttt{XOR}) & \texttt{111001} $= \{1, 2, 3, 6\}$ \\
        $a \cup b$ & $a + b + ab$ & \texttt{a | b} (\texttt{OR}) & \texttt{111011} $= \{1, 2, 3, 5, 6\}$ \\
        $I_n \setminus a$ & $1 + a$ & \texttt{\textasciitilde a} (\texttt{NOT}) & \texttt{100101} $= \{1, 4, 6\}$ \\
    \end{tabular}
    \end{center}
    È il modo in cui i programmi memorizzano insiemi piccoli in un solo numero intero (le «maschere di bit»): ad esempio i permessi
    \texttt{rwx} di un file sono una parte di $\{\text{lettura}, \text{scrittura}, \text{esecuzione}\}$, cioè una stringa di $3$ bit.
\end{spiegazione}

\begin{concetto}
\begin{Corollario}{Gli anelli booleani finiti sono le stringhe di bit}{booleani-z2n}
    \begin{prerequisiti}
        \dip{Corollario}{cor:conseguenze-stone}
        \dip{Teorema}{teo:caratteristiche-z2n}
    \end{prerequisiti}
    Ogni anello booleano finito è isomorfo a $(\mbZ_2^n, +, \cdot)$ per un unico $n \geq 1$, quello per cui il suo ordine è $2^n$.
\end{Corollario}
\begin{dimostrazione}
    Se l'anello ha ordine $2^n$ (Corollario~\ref{cor:conseguenze-stone}), ha lo stesso ordine di $\mbZ_2^n$, che è un anello booleano: per il Corollario~\ref{cor:conseguenze-stone} i due sono
    isomorfi. L'unicità di $n$ segue dal fatto che $2^n$ determina $n$.
\end{dimostrazione}
\end{concetto}

\begin{spiegazione}
    Ricapitolando, per ogni $n \geq 1$ c'è, a meno di isomorfismi, un solo anello booleano con $2^n$ elementi, e lo possiamo vedere in quattro modi
    equivalenti:
    \[ (\mcP(I_n), \subseteq) \quad\longleftrightarrow\quad (\mcP(I_n), \cup, \cap, \varnothing, I_n, {}') \quad\longleftrightarrow\quad (\mcP(I_n), \Delta, \cap) \quad\longleftrightarrow\quad (\mbZ_2^n, +, \cdot). \]
    Con $n = 4$, ad esempio, i divisori di $210$, le parti di $\{2, 3, 5, 7\}$ e le stringhe di $4$ bit sono lo stesso anello booleano: il divisore $15 = 3 \cdot 5$ corrisponde
    a $\{3, 5\}$ e alla stringa \texttt{0110} (i bit dicono se compaiono $2$, $3$, $5$, $7$).
\end{spiegazione}
